\environment env_ConfigGlobal
\environment env_Config

\startcomponent AppendixExerciseAnswers
\startchapter[title={Answers to lesson exercises}]

\startsubject[title={Lesson \in{}[sec:exercises1] answers}]
\startitemize[n,packed]
\item \JPN{わたしがねこだ} (watashi-ga neko da)
\item \JPN{わたしがおどる} (watashi-ga odoru)
\item \JPN{ねこがあるく} (neko-ga aruku)
\item \JPN{とりがうたう} (tori-ga utau)
\item \JPN{ねこがしろい} (neko-ga shiroi)
\item \JPN{とりが\ruby{小}{ちい}さい} (tori-ga chiisai)
\stopitemize
\stopsubject

\startsubject[title={Lesson \in{}[sec:exercises2] answers}]
\startitemize[n]

\startitem
\JPN{ありすだ} (Arisu da)\crlf
=\JPN{\varnothing{}がありすだ} (=\varnothing{}ga Arisu da)\crlf
=\JPN{わたしがありすだ} (=watashi-ga Arisu da)
\stopitem

\startitem
\JPN{はれだ} (hare da)\crlf
=\JPN{\varnothing{}がはれだ} (=\varnothing{}ga hare da)\crlf
=\JPN{\ruby{天}{てん}\ruby{気}{き}がはれだ} (=tenki-ga hare da)\crlf
or =\JPN{\ruby{今日}{きょう}がはれだ} (=kyou-ga hare da)
\stopitem

\startitem
\JPN{おちゃを\ruby{飲}{の}む} (ocha-wo nomu)\crlf
=\JPN{\varnothing{}がおちゃを\ruby{飲}{の}む} (=\varnothing{}ga ocha-wo nomu)\crlf
=\JPN{わたしがおちゃを\ruby{飲}{の}む} (=watashi-ga ocha-wo nomu)
\stopitem

\startitem
\JPN{ほわいとでーだ} (howaito dē da)\crlf
=\JPN{\varnothing{}がほわいとでーだ} (=\varnothing{}ga howaito dē da)\crlf
=\JPN{\ruby{今日}{きょう}がほわいとでーだ} (=kyou-ga howaito dē da)
\stopitem

\startitem
\JPN{うたをうたう} (uta-wo utau)\crlf
=\JPN{\varnothing{}がうたをうたう} (=\varnothing{}ga uta-wo utau)\crlf
=\JPN{わたしがうたをうたう} (=watashi-ga uta-wo utau)
\stopitem

\stopitemize
\stopsubject

\startsubject[title={Lesson \in{}[sec:exercises3] Answers}]

The comma is just to show that the は-marked topic-flag is not a part of the
logical clause. If you don't have it, that's all right. If you have the に- and
を-marked nouns in the opposite order shown, that's OK too. What matters is having
the right particles attached to the right nouns.

\startitemize[n]

\startitem
\JPN{わたしは、\Sakura{}だ} (watashi-wa, Sakura da)\crlf
=\JPN{わたしは、\varnothing{}が\Sakura{}だ} (=watashi-wa, \varnothing{}ga Sakura da)\crlf
=\JPN{わたしは、わたしが\Sakura{}だ} (=watashi-wa, watashi-ga Sakura da)\crlf
=As for me, I am Sakura
\stopitem

\startitem
\JPN{\Sakura{}は、とてもうつくしい} (Sakura-wa, totemo utsukushii)\crlf
=\JPN{\Sakura{}は、\varnothing{}がとてもうつくしい} (=Sakura-wa, \varnothing{}ga totemo utsukushii)\crlf
=\JPN{\Sakura{}は、\Sakura{}がとてもうつくしい} (=Sakura-wa, Sakura-ga totemo utsukushii)\crlf
=As for Sakura, she is very beautiful
\stopitem

\startitem
\JPN{わたしは、\Sakura{}にてがみをおくる} (watashi-wa, Sakura-ni tegami-wo okuru)\crlf
=\JPN{わたしは、\varnothing{}が\Sakura{}にてがみをおくる} (=watashi-wa, \varnothing{}ga Sakura-ni tegami-wo okuru)\crlf
=\JPN{わたしは、わたしが\Sakura{}にてがみをおくる} (=watashi-wa, watashi-ga Sakura-ni tegami-wo okuru)\crlf
=As for me, I send a letter to Sakura
\stopitem

\startitem
\JPN{てがみは、\Sakura{}におくる} tegami-wa, Sakura-ni okuru)\crlf
=\JPN{てがみは、\varnothing{}が\varnothing{}を\Sakura{}におくる} (=tegami-wa, \varnothing{}ga \varnothing{}wo Sakura-ni okuru)\crlf
=\JPN{てがみは、わたしがてがみを\Sakura{}におくる} (=tegami-wa, watashi-ga tegami-wo Sakura-ni okuru)\crlf
=As for the letter, I send it to Sakura\crlf
\EditorsNote{Did this answer surprise you? The reason we can have two zero-pronouns in this sentence is because the the subject ({\em I}) and the topic ({\em it}: the letter) are both implied. The default value of \varnothing{}が is \JPN{わたしが}, leaving the topic as the default value for \varnothing{}を \quote{it}. Think of this casual English sentence where {\em I} is invisible, but still there: \quotation{Speaking of that letter, gotta send it to Sakura.}}
\stopitem

\startitem
\JPN{\Sakura{}は、てがみをおくる} (Sakura-wa, tegami-wo okuru)\crlf
=\JPN{\Sakura{}は、\varnothing{}が\varnothing{}にてがみをおくる} (=Sakura-wa, \varnothing{}ga \varnothing{}ni tegami-wo okuru)\crlf
=\JPN{\Sakura{}は、わたしが\Sakura{}にてがみをおくる} (=Sakura-wa, watashi-ga Sakura-ni tegami-wo okuru)\crlf
=As for Sakura, I send her a/the letter\crlf
\EditorsNote{Again, we can have two zero-pronouns in the sentence because the implied subject is \JPN{わたしが}, leaving the topic {\em Sakura} as the default value for \varnothing{}に \quote{her}.}
\stopitem

\startitem
\JPN{\Sakura{}は、わたしにてがみをおくる} (Sakura-wa, watashi-ni tegami-wo okuru)\crlf
=\JPN{\Sakura{}は、\varnothing{}がわたしにてがみをおくる} (=Sakura-wa, \varnothing{}ga watashi-ni tegami-wo okuru)\crlf
=\JPN{\Sakura{}は、\Sakura{}がわたしにてがみをおくる} (=Sakura-wa, Sakura-ga watashi-ni tegami-wo okuru)\crlf
=As for Sakura, she sends me a/the letter
\stopitem

\stopitemize
\stopsubject

\stopchapter
\stopcomponent
